% Sources: L01-L03, L21-L25, L27-L38.
\boardsection{4}{假设检验}
\[
H_0:\theta\in\Theta_0,\qquad H_1:\theta\in\Theta_1,\qquad
\Theta_0\cap\Theta_1=\varnothing,\quad\Theta_0\cup\Theta_1=\Theta.
\]
例如$H_0:\theta\leq0.05$，$H_1:\theta>0.05$。
\[
\alpha_i=P(H_i\mid x)=\int_{\Theta_i}\pi(\theta\mid x)\,d\theta,
\quad i=0,1,\qquad\alpha_0+\alpha_1=1.
\]
例如后验为$\Be(a,b)$，
\[
\alpha_0=F(0.05)=0.8,\qquad\alpha_1=0.2.
\]
若误判代价$C_{01}=10,\ C_{10}=1$，
\[
R(\text{选择}H_0\mid x)=10\times0.2=2,\qquad
R(\text{选择}H_1\mid x)=1\times0.8=0.8.
\]
选择$H_1$。

频率检验：拒绝域$W$，显著性水平$\alpha$，
\[
\sup_{\theta\in\Theta_0}P_\theta(X\in W)\leq\alpha.
\]

\block{贝叶斯因子}
记两个假设的先验与后验概率为
\[
\pi_i=P(H_i),\quad\alpha_i=P(H_i\mid x),\quad i=0,1.
\]
且$\pi_0,\pi_1$均大于0。支持$H_0$相对于$H_1$的Bayes因子定义为
\[
\BF_{01}(x)=\frac{\alpha_0/\alpha_1}{\pi_0/\pi_1},
\qquad
\underbrace{\frac{\alpha_0}{\alpha_1}}_{\text{后验赔率}}
=\underbrace{\frac{\pi_0}{\pi_1}}_{\text{先验赔率}}\BF_{01}(x).
\]
贝叶斯因子等于边际似然比。记假设$H_i$下data的边际密度为
\[
m_i(x)=p(x\mid H_i).
\]
由贝叶斯公式，
\[
\alpha_i=\frac{\pi_i m_i(x)}{\pi_0m_0(x)+\pi_1m_1(x)},\quad i=0,1.
\]
\[
\frac{\alpha_0}{\alpha_1}
=\frac{\pi_0m_0(x)}{\pi_1m_1(x)},\qquad
\BF_{01}(x)=\frac{m_0(x)}{m_1(x)}.
\]
对于复合假设，
\[
m_i(x)=\int_{\Theta_i}p(x\mid\theta)\pi(\theta\mid H_i)\,d\theta.
\]
\[
g_i(\theta)=\frac{\pi(\theta)\ind_{\Theta_i}(\theta)}{\pi_i},
\quad
\pi(\theta)=\pi_0g_0(\theta)+\pi_1g_1(\theta),
\]
\[
\BF_{01}(x)=
\frac{\int_{\Theta_0}p(x\mid\theta)g_0(\theta)\,d\theta}
{\int_{\Theta_1}p(x\mid\theta)g_1(\theta)\,d\theta}.
\]

\block{例2.4.1}
\[
X\mid\theta\sim\Bin(n,\theta),\quad\theta\sim U(0,1),
\qquad H_0:\theta\leq\tfrac12,\quad H_1:\theta>\tfrac12.
\]
\[
\theta\mid x\sim\Be(x+1,n-x+1).
\]
\[
\alpha_0=P(\theta\leq\tfrac12\mid x)
=\frac{(n+1)!}{x!(n-x)!}
\int_0^{1/2}\theta^x(1-\theta)^{n-x}\,d\theta,\quad\alpha_1=1-\alpha_0.
\]
令$m=n-x$，
\[
I_{x,m}=\int_0^{1/2}\theta^x(1-\theta)^m\,d\theta.
\]
分部积分：
\begin{align*}
I_{x,m}
&=\left[\frac{\theta^{x+1}(1-\theta)^m}{x+1}\right]_0^{1/2}
+\frac m{x+1}\int_0^{1/2}\theta^{x+1}(1-\theta)^{m-1}\,d\theta\\
&=\frac{2^{-(n+1)}}{x+1}+\frac m{x+1}I_{x+1,m-1}.
\end{align*}
\[
I_{n,0}=\frac{2^{-(n+1)}}{n+1}.
\]
\begin{align*}
I_{x,m}
=\frac{2^{-(n+1)}}{x+1}
\left[1+\frac m{x+2}
+\frac{m(m-1)}{(x+2)(x+3)}
+\cdots+\frac{m!(x+1)!}{(n+1)!}\right].
\end{align*}
另一种展开：
\[
(1-\theta)^m=\sum_{k=0}^m(-1)^k\binom mk\theta^k,
\quad
I_{x,m}=\sum_{k=0}^m(-1)^k\binom mk
\frac{2^{-(x+k+1)}}{x+k+1}.
\]
\[
\alpha_0=\frac{(n+1)!}{x!(n-x)!}
\sum_{k=0}^{n-x}(-1)^k\binom{n-x}{k}
\frac{2^{-(x+k+1)}}{x+k+1}.
\]
取$n=5,x=1$，则
\begin{align*}
\alpha_0
&=\frac{6!}{1!4!}
\left(\frac18-\frac16+\frac3{32}-\frac1{40}+\frac1{384}\right)\\
&=30\times\frac{57}{1920}=\frac{57}{64},\qquad
\alpha_1=\frac7{64}.
\end{align*}
Beta分布的转化性质。若$U\sim\Be(a,b)$，$a,b\in\{1,2,\ldots\}$，
对$0\leq t\leq1$，
\[
P(U\leq t)=\sum_{j=a}^{a+b-1}\binom{a+b-1}{j}
t^j(1-t)^{a+b-1-j}.
\]
若$Y\sim\Bin(a+b-1,t)$，则$P(U\leq t)=P(Y\geq a)$。
代入$a=x+1,\ b=n-x+1$，
\[
\alpha_0=P\{\Bin(n+1,\tfrac12)\geq x+1\}
=\frac1{2^{n+1}}\sum_{j=x+1}^{n+1}\binom{n+1}{j}.
\]
如$n=5,x=1$，
\[
\alpha_0=1-P(Y=0)-P(Y=1)
=1-\frac{\binom60+\binom61}{64}=\frac{57}{64}.
\]
\begin{center}
\begin{tabular}{c c c c}\toprule
$x$&$P(H_0\mid x)$&$P(H_1\mid x)$&后验赔率\\\midrule
0&$63/64\approx0.984$&$1/64$&$63:1$\\
1&$57/64\approx0.891$&$7/64$&$8.14:1$\\
2&$42/64\approx0.656$&$22/64$&$1.91:1$\\
3&$22/64$&$42/64$&$0.52$\\
4&$7/64$&$57/64$&$0.12$\\
5&$1/64$&$63/64$&$0.016$\\\bottomrule
\end{tabular}
\end{center}
$x=0,1,2$选择$H_0$；$x=3,4,5$选择$H_1$。

注：对称性，$\alpha_0(x)=\alpha_1(n-x)$，即交换成功与失败次数，就交换了对应的假设支持。

\block{简单原假设与简单备择假设}
\[
H_0:\theta=\theta_0,\qquad H_1:\theta=\theta_1.
\]
\[
m(x)=\pi_0p(x\mid\theta_0)+\pi_1p(x\mid\theta_1),
\quad
\alpha_0=\frac{\pi_0p(x\mid\theta_0)}{m(x)},\quad
\alpha_1=\frac{\pi_1p(x\mid\theta_1)}{m(x)}.
\]
\[
\frac{\alpha_0}{\alpha_1}
=\frac{\pi_0p(x\mid\theta_0)}{\pi_1p(x\mid\theta_1)},\qquad
\BF_{01}(x)=\frac{\alpha_0/\alpha_1}{\pi_0/\pi_1}
=\frac{p(x\mid\theta_0)}{p(x\mid\theta_1)}.
\]
$\BF_{01}(x)>1$，证据向$H_0$倾斜；
$\BF_{01}(x)<1$，证据向$H_1$倾斜。

在两种错误损失相同、正确决策损失为0时，选择后验概率较大的假设。
选择$H_1$当且仅当$\alpha_1>\alpha_0$，等价于
\[
\frac{p(x\mid\theta_1)}{p(x\mid\theta_0)}>\frac{\pi_0}{\pi_1}.
\]
N-P引理：
\[
\frac{p(x\mid\theta_1)}{p(x\mid\theta_0)}>C\qquad(\text{阈值}).
\]
例2.10中$P(H_0)=P(H_1)=1/2$，故$\BF_{01}$恰好等于后验赔率。
$n=5,x=1$时，
\[
\BF_{01}(x)=\frac{57}{7}\approx8.14,\qquad
P(H_0\mid x)=\frac{8.14}{1+8.14}\approx0.891.
\]

\block{例2.11}
\[
X_1,\ldots,X_n\overset{\mathrm{iid}}{\sim}N(\theta,1),\quad
\theta\in\{0,1\},\qquad H_0:\theta=0,\quad H_1:\theta=1.
\]
\[
\sum_i(x_i-\theta)^2=\sum_i(x_i-\bar x)^2+n(\bar x-\theta)^2.
\]
略写似然函数：
\[
L(\theta;x)=h(x)\exp\left\{-\frac n2(\bar x-\theta)^2\right\}.
\]
由因子分解定理，$\bar X$为充分统计量，比较两个假设时$h(x)$约去，可直接用
$\bar X\mid\theta\sim N(\theta,1/n)$。
\begin{align*}
\BF_{01}(x)
&=\frac{p(x\mid\theta=0)}{p(x\mid\theta=1)}\\
&=\exp\left\{-\frac n2\bar x^2+\frac n2(\bar x-1)^2\right\}\\
&=\exp\left\{-\frac n2(2\bar x-1)\right\}.
\end{align*}
\[
\bar x<\tfrac12:\ \BF_{01}>1,\text{支持}H_0;\quad
\bar x>\tfrac12:\ \BF_{01}<1,\text{支持}H_1;\quad
\bar x=\tfrac12:\ \BF_{01}=1.
\]
若$n=10,\bar x=2$，有
\[
\BF_{01}=e^{-10(4-1)/2}=e^{-15}\approx3.06\times10^{-7},
\quad
\BF_{10}=e^{15}\approx3.3\times10^6.
\]
①若两种假设的先验概率相等，
\[
P(H_0\mid x)=\frac{\BF_{01}}{1+\BF_{01}}\approx\BF_{01},
\quad P(H_1\mid x)\approx1-\BF_{01}\approx0.999999694.
\]
②由
\[
\frac{P(H_0\mid x)}{P(H_1\mid x)}
=\BF_{01}\frac{\pi_0}{\pi_1},
\]
欲选择$H_0$，必须满足
\[
\BF_{01}\frac{\pi_0}{\pi_1}>1,\qquad
e^{-15}\frac{\pi_0}{\pi_1}>1,\qquad
\frac{\pi_0}{\pi_1}>e^{15}\approx330\text{万}.
\]
即使$\pi_0/\pi_1=10000$，后验赔率$\approx3.06\times10^{-3}$，
对应$P(H_0\mid x)\approx3.04\times10^{-3}$，仍选择$H_1$。
\[
\frac{2-1}{1/\sqrt{10}}=\sqrt{10}\approx3.2.
\]

\block{注：广义似然比与后验赔率}
\[
\lambda(x)=\frac{\sup_{\theta\in\Theta_0}p(x\mid\theta)}
{\sup_{\theta\in\Theta_0\cup\Theta_1}p(x\mid\theta)}.
\]
\[
\text{后验赔率}=\text{先验赔率}\times
\frac{H_0\text{下的边际似然}}{H_1\text{下的边际似然}}.
\]

\block{例2.4.3}
\[
X_1,\ldots,X_n\overset{\mathrm{iid}}{\sim}N(\theta,1),\qquad
\theta\sim N(\mu,2).
\]
\[
\pi(\theta\mid x)\propto
\exp\left\{-\frac12\left[n(\theta-\bar x)^2+\frac12(\theta-\mu)^2\right]\right\}.
\]
配方：
\[
\theta\mid x\sim N(\mu_n,S_n^2),\qquad
\mu_n=\frac{n\bar x+\mu/2}{n+1/2},\quad S_n^2=\frac1{n+1/2}.
\]
取$n=10,\bar x=1.5,\mu=0.5$，
\[
\mu_n=\frac{15+0.25}{10.5}=1.452381,\qquad
S_n=\sqrt{\frac1{10.5}}=0.308607.
\]
\[
H_0=\{\theta\leq1\},\qquad
\alpha_0=P(H_0\mid x)=\Phi\left(\frac{1-\mu_n}{S_n}\right).
\]
\[
\frac{1-\mu_n}{S_n}\approx-1.46585,\qquad
\alpha_0\approx0.07134,\quad\alpha_1\approx0.92866,\quad
\frac{\alpha_0}{\alpha_1}\approx0.07682.
\]
\[
\pi_0=P(\theta\leq1)
=\Phi\left(\frac{1-0.5}{\sqrt2}\right)\approx0.63816,\qquad
\pi_1\approx0.36184,\quad\frac{\pi_0}{\pi_1}\approx1.76366.
\]
\[
\BF_{01}\approx\frac{0.07682}{1.76366}\approx0.04356,\qquad
\BF_{10}\approx22.96.
\]

\block{3．简单原假设与复合备择假设}
\[
H_0:\theta=\theta_0,\qquad H_1:\theta\ne\theta_0.
\]
设$P(H_0)=\pi_0$，$P(H_1)=\pi_1=1-\pi_0$。
在$H_1$下给定$\theta$的密度$g_1(\theta)$，则
\[
\pi(d\theta)=\pi_0\delta_{\theta_0}(d\theta)+\pi_1g_1(\theta)\,d\theta,
\qquad
\delta_{\theta_0}(A)=\ind_{\{\theta_0\in A\}}.
\]
在$H_0$下$m_0(x)=p(x\mid\theta_0)$；在$H_1$下
\[
m_1(x)=\int_{\theta\ne\theta_0}p(x\mid\theta)g_1(\theta)\,d\theta.
\]
\[
m(x)=\pi_0p(x\mid\theta_0)+\pi_1m_1(x).
\]
\[
\alpha_0=\frac{\pi_0p(x\mid\theta_0)}
{\pi_0p(x\mid\theta_0)+\pi_1m_1(x)},\qquad
\alpha_1=\frac{\pi_1m_1(x)}
{\pi_0p(x\mid\theta_0)+\pi_1m_1(x)}.
\]
\[
\frac{\alpha_0}{\alpha_1}=\frac{\pi_0p(x\mid\theta_0)}{\pi_1m_1(x)},\quad
\BF_{01}(x)=\frac{p(x\mid\theta_0)}{m_1(x)},\quad
P(H_0\mid x)=
\left[1+\frac{1-\pi_0}{\pi_0}\frac1{\BF_{01}(x)}\right]^{-1}.
\]
如$\pi_0=1/2$，
\[
P(H_0\mid x)=\frac{\BF_{01}}{1+\BF_{01}}.
\]

\block{例2.4.4}
\[
X\mid\theta\sim\Bin(n,\theta),\qquad
\pi(d\theta)=\pi_0\delta_{1/2}(d\theta)
+(1-\pi_0)\ind_{(0,1)}(\theta)\,d\theta.
\]
\[
P(\theta=1/2)=\pi_0>0.
\]
\begin{align*}
m_1(x)&=\int_0^1\binom nx\theta^x(1-\theta)^{n-x}\,d\theta\\
&=\binom nx B(x+1,n-x+1)=\frac1{n+1},\quad x=0,\ldots,n,\\
m_0(x)&=P(X=x\mid\theta=1/2)=\binom nx2^{-n}.
\end{align*}
\[
\BF_{01}(x)=\frac{\binom nx2^{-n}}{1/(n+1)}
=\frac{(n+1)!}{2^n x!(n-x)!}.
\]
取$n=5,x=3$，
\[
m_0=\frac{\binom53}{2^5}=\frac5{16},\qquad
m_1=\frac16,\qquad\BF_{01}=\frac{15}{8}=1.875.
\]
若$\pi_0=\pi_1=1/2$，
\[
P(H_0\mid x)=\frac{15}{23}\approx0.6522,\qquad
P(H_1\mid x)=\frac8{23}\approx0.3478.
\]
选择$H_0$。
